\documentclass[11pt,a4paper]{article}

% --- Encoding ---
\usepackage[utf8]{inputenc}
\usepackage[T1]{fontenc}
\usepackage[spanish,es-noshorthands]{babel}
\usepackage{lmodern}

% --- Layout ---
\usepackage[margin=2.5cm]{geometry}
\usepackage{parskip}
\setlength{\parindent}{0pt}
\setlength{\parskip}{0.5em plus 0.2em minus 0.1em}

% --- Tables and graphics ---
\usepackage{booktabs}
\usepackage{array}
\usepackage{enumitem}
\usepackage{xcolor}
\usepackage{graphicx}
\usepackage{tikz}
\usetikzlibrary{positioning,shapes.geometric,arrows.meta,fit,backgrounds}
\usepackage{caption}
\usepackage{amsmath,amssymb}

% --- Code listings ---
\usepackage{listings}
\definecolor{codegray}{rgb}{0.96,0.96,0.96}
\definecolor{codeborder}{rgb}{0.85,0.85,0.85}
\definecolor{codecomment}{rgb}{0.30,0.55,0.30}
\definecolor{codekeyword}{rgb}{0.10,0.30,0.70}
\definecolor{codestring}{rgb}{0.65,0.15,0.10}

\lstdefinestyle{cstyle}{
    language=C,
    backgroundcolor=\color{codegray},
    basicstyle=\ttfamily\footnotesize,
    breaklines=true,
    frame=single,
    rulecolor=\color{codeborder},
    showstringspaces=false,
    keywordstyle=\color{codekeyword}\bfseries,
    commentstyle=\color{codecomment}\itshape,
    stringstyle=\color{codestring},
    numbers=left,
    numberstyle=\tiny\color{gray},
    numbersep=8pt,
    aboveskip=1em,
    belowskip=1em,
}

\lstdefinestyle{shellstyle}{
    language=bash,
    backgroundcolor=\color{codegray},
    basicstyle=\ttfamily\small,
    breaklines=true,
    frame=single,
    rulecolor=\color{codeborder},
    showstringspaces=false,
    keywordstyle=\color{codekeyword}\bfseries,
    commentstyle=\color{codecomment}\itshape,
    stringstyle=\color{codestring},
    aboveskip=1em,
    belowskip=1em,
    columns=fullflexible,
}

\lstdefinestyle{textstyle}{
    backgroundcolor=\color{codegray},
    basicstyle=\ttfamily\footnotesize,
    breaklines=true,
    frame=single,
    rulecolor=\color{codeborder},
    showstringspaces=false,
    aboveskip=1em,
    belowskip=1em,
    columns=fullflexible,
}

\lstset{style=shellstyle}

% --- Definition / Idea / Important boxes ---
\usepackage[most]{tcolorbox}
\newtcolorbox{definicion}[1][]{
    colback=blue!5!white,
    colframe=blue!50!black,
    fonttitle=\bfseries,
    title=Definición,
    sharp corners=southwest,
    rounded corners=northwest,
    arc=2pt,
    boxrule=0.6pt,
    left=8pt,right=8pt,top=4pt,bottom=4pt,
    #1
}
\newtcolorbox{idea}[1][]{
    colback=orange!5!white,
    colframe=orange!60!black,
    fonttitle=\bfseries,
    title=Idea intuitiva,
    sharp corners=southwest,
    rounded corners=northwest,
    arc=2pt,
    boxrule=0.6pt,
    left=8pt,right=8pt,top=4pt,bottom=4pt,
    #1
}
\newtcolorbox{importante}[1][]{
    colback=red!5!white,
    colframe=red!50!black,
    fonttitle=\bfseries,
    title=Importante,
    sharp corners=southwest,
    rounded corners=northwest,
    arc=2pt,
    boxrule=0.6pt,
    left=8pt,right=8pt,top=4pt,bottom=4pt,
    #1
}

% --- Hyperlinks ---
\usepackage[colorlinks=true,
            linkcolor=blue!60!black,
            urlcolor=blue!50!black,
            citecolor=blue!60!black]{hyperref}

% --- Title ---
\title{
    \vspace{-2.5em}
    \textbf{Resumen teórico --- Sistemas Operativos} \\[0.3em]
    \large Notas del curso 5 y 6: \emph{Procesos, threads, concurrencia e IPC}
}
\author{
    Tomás Rodeghiero \\
    \small Sistemas Operativos --- Universidad Nacional de Río Cuarto
}
\date{2026}

\begin{document}

\maketitle

\begin{abstract}
\noindent
Este documento es un resumen teórico de los capítulos 5
(\emph{Procesos y threads}) y 6 (\emph{Concurrencia e IPC}) de las
\emph{Notas del curso} de Sistemas Operativos (Marcelo Arroyo, UNRC).
Se cubren las nociones de proceso y de thread, los estados por los que
puede pasar una tarea, las llamadas al sistema asociadas a su gestión,
los modos de ejecución del procesador, el mecanismo de interrupciones
y traps, los cambios de contexto, la concurrencia y sus problemas
clásicos (condiciones de carrera, secciones críticas), las primitivas
de sincronización (locks, spinlocks, sleep locks, semáforos, monitores,
variables de condición), los problemas asociados al uso de locks
(\emph{deadlock}, \emph{livelock}), los modelos de comunicación por
mensajes y los mecanismos IPC ofrecidos por UNIX. El objetivo es servir
como material de estudio para preparar el Práctico~3 y el primer
parcial: cada concepto se introduce con intuición, se formaliza cuando
hace falta y se conecta con los demás para que el lector pueda razonar
sobre el modelo completo del SO.
\end{abstract}

\tableofcontents
\bigskip

% =====================================================================
\section{Introducción}
% =====================================================================

Cuando escribimos un programa pensamos en una secuencia única de
instrucciones que se ejecuta de principio a fin. La realidad de un
sistema operativo moderno es muy distinta: en una computadora hay
\emph{decenas} de programas corriendo a la vez --- el navegador, el
editor de texto, el reloj de la barra de tareas, los servicios del
sistema, el kernel mismo --- y la CPU es un recurso limitado. Una de
las funciones más importantes del SO es darle a cada programa la
\emph{ilusión} de tener una computadora dedicada, mientras se ocupa
internamente de repartir la CPU entre todos.

Para hacer esto posible el SO necesita una serie de abstracciones:

\begin{itemize}[leftmargin=*]
    \item un \textbf{proceso} representa cada programa en ejecución;
    \item un \textbf{thread} es la unidad más chica que el scheduler
        del SO puede planificar;
    \item un \textbf{estado} describe en qué momento de su ciclo de
        vida está cada tarea;
    \item un \textbf{cambio de contexto} es la operación con la que el
        SO le saca la CPU a una tarea y se la da a otra;
    \item los mecanismos de \textbf{sincronización} y los de
        \textbf{IPC} permiten que las tareas cooperen sin pisarse.
\end{itemize}

\begin{idea}
La intuición es: el SO simula varias computadoras virtuales sobre una
única física. Para que funcione, tiene que poder \emph{guardar el
trabajo en curso} de una tarea, cambiar a otra, y luego retomar la
primera exactamente donde la dejó.
\end{idea}

% =====================================================================
\section{Procesos y threads}
% =====================================================================

\subsection{Definiciones básicas}

\begin{definicion}
Un \textbf{proceso} es una instancia de un programa en ejecución. Un
\textbf{programa} es un archivo ejecutable que reside en algún
dispositivo de almacenamiento. Un proceso tiene al menos un thread.
\end{definicion}

\begin{definicion}
Un \textbf{thread} (o hilo) es la mínima secuencia de instrucciones
que el SO puede planificar para ejecutarse en una CPU.
\end{definicion}

Un programa \emph{en disco} es estático: no consume CPU ni memoria del
sistema. Solo cuando es \emph{cargado} en memoria y se le asigna un
thread de ejecución se convierte en un \emph{proceso}. Esa distinción
es fundamental: del mismo programa pueden existir muchos procesos
ejecutando simultáneamente, cada uno con su propio estado.

Un proceso ejecuta el código en \emph{modo usuario} (el modo de menor
privilegio de la CPU) y solo puede acceder a su \emph{propio espacio
de memoria}. Esta protección la implementa el SO con la ayuda de la
MMU del hardware: si una instrucción intenta tocar una dirección fuera
del espacio del proceso, la CPU dispara una excepción.

\subsection{Tareas: la abstracción común}

Es conveniente abstraer las diferencias técnicas entre procesos y
threads y llamar a ambos \textbf{tareas} (\emph{tasks}). Tanto un
proceso como un thread son tareas que el scheduler conoce, planifica
y representa con un descriptor.

Un \emph{proceso multithreaded} crea varios threads dentro de él
mismo. Cada thread tiene su propia pila, pero todos comparten el
espacio de memoria del proceso (el código, los datos globales, el
heap y los descriptores de archivos abiertos).

\begin{center}
\begin{tikzpicture}[node distance=0pt, every node/.style={draw, minimum width=4.5cm, minimum height=0.5cm, font=\ttfamily\small}]
    \node[fill=blue!8!white] (code)  {code (instructions)};
    \node[fill=blue!8!white,below=of code]  (data)  {global data};
    \node[fill=gray!10,below=of data] (free) {... free space ...};
    \node[fill=green!10,below=of free] (st1) {stack thread 1};
    \node[below=of st1] (dots) {...};
    \node[fill=green!10,below=of dots] (stm) {stack main thread};
\end{tikzpicture}
\\
\small\emph{Figura 1: Imagen en memoria de un proceso multithreaded.}
\end{center}

\subsection{Atributos de un proceso}

El SO mantiene, por cada proceso, un descriptor (a veces llamado
\emph{Process Control Block}, PCB) con al menos:

\begin{itemize}[leftmargin=*]
    \item un identificador único (\texttt{pid});
    \item su estado actual;
    \item si está dormido, la razón o recurso por el cual lo está;
    \item su \emph{mapa de memoria} accesible en modo usuario:
    \begin{itemize}
        \item segmento de código (\emph{text});
        \item datos globales inicializados (\emph{data}) y sin
            inicializar (\emph{bss});
        \item heap (memoria dinámica, manejada con \texttt{sbrk} o
            \texttt{malloc});
        \item stacks de cada thread.
    \end{itemize}
    \item al menos un thread de ejecución, representado por su
        \emph{stack en modo usuario} y su \emph{stack en modo kernel};
    \item el conjunto de recursos adquiridos (archivos abiertos,
        pipes, sockets, semáforos, áreas de memoria compartida,
        manejadores de señales, timers);
    \item el comando o programa del cual es instancia.
\end{itemize}

El \emph{stack en modo kernel} merece atención especial: en él se
guarda el \emph{trapframe} (los valores de los registros del proceso
al momento de una interrupción), se usa para las invocaciones a
funciones del kernel mientras se ejecuta en el contexto del proceso
y, finalmente, se guarda el \emph{contexto} (registros caller-saved)
del thread cuando ocurre un cambio de contexto.

\subsection{Layout virtual de memoria}

\begin{center}
\begin{tikzpicture}[every node/.style={draw, minimum width=3.5cm, minimum height=0.55cm, font=\ttfamily\small}]
    \node[fill=blue!8!white] (text) {text (code)};
    \node[fill=blue!8!white,below=0pt of text] (sd)   {static data};
    \node[fill=gray!10, below=0pt of sd, minimum height=1.4cm] (free) {free space};
    \node[fill=green!10,below=0pt of free, minimum height=0.8cm] (stk) {stack};
    \node[draw=none,right=0.4cm of sd]   (brk) {\small program break (brk)};
    \node[draw=none,right=0.4cm of stk]  (args) {\small main arguments};
    \node[draw=none, left=0.4cm of text] (start) {\small start};
    \node[draw=none, left=0.4cm of stk]  (end)   {\small end};
\end{tikzpicture}
\\
\small\emph{Figura 2: Layout virtual de la memoria de un proceso.}
\end{center}

El \emph{program break} (\texttt{brk}) marca el final del segmento de
datos. La syscall \texttt{sbrk(n)} permite extender o reducir el área
de datos asignando o liberando memoria al proceso (lo que usa
\texttt{malloc()} internamente). Las direcciones \texttt{start} y
\texttt{end} delimitan el espacio virtual del proceso; en muchos
sistemas \texttt{start = 0}.

\subsection{Registros del thread y el calling convention}

El estado de un thread en un punto cualquiera de su ejecución está
representado por los valores de los registros de la CPU. Tres son
particularmente importantes:

\begin{itemize}[leftmargin=*]
    \item \textbf{Program counter} (\texttt{pc}): dirección de la
        próxima instrucción a ejecutar.
    \item \textbf{Stack pointer} (\texttt{sp}): tope de la pila.
    \item \textbf{Frame (base) pointer} (\texttt{fp}): dirección base
        del registro de activación de la función actual.
\end{itemize}

La pila contiene \emph{registros de activación} (\emph{activation
records} o \emph{stack frames}) para controlar la secuencia de
invocación a funciones y sus retornos. Cada registro de activación
contiene la dirección de retorno, el frame pointer salvado, los
argumentos y las variables locales de la función. Esta organización
es la que permite que las llamadas recursivas funcionen correctamente
y que las variables locales se creen y destruyan automáticamente.

El \textbf{calling convention} de la ABI de la plataforma especifica
cómo se pasan los argumentos (en registros, en pila, o en una
combinación) y cómo se devuelve el valor de retorno. En RISC-V, por
ejemplo:

\begin{itemize}[leftmargin=*,nosep]
    \item \texttt{x1}/\texttt{ra}: dirección de retorno;
    \item \texttt{x2}/\texttt{sp}: stack pointer;
    \item \texttt{x8}/\texttt{fp}: frame pointer;
    \item \texttt{x10}--\texttt{x17}/\texttt{a0}--\texttt{a7}:
        argumentos (los excedentes se apilan);
    \item \texttt{a0}/\texttt{a1}: valor de retorno;
    \item \texttt{s0}--\texttt{s11}: \emph{saved registers} que la
        función invocada debe preservar.
\end{itemize}

\subsection{Llamadas al sistema para gestión de procesos}

Las syscalls clásicas de POSIX para administrar procesos son:

\begin{center}
\small
\begin{tabular}{@{}lp{8.5cm}@{}}
\toprule
Syscall & Qué hace \\
\midrule
\texttt{fork()} & Crea un proceso hijo, copia exacta del padre. El hijo hereda los descriptores de archivos abiertos. \\
\texttt{exec(path, args)} & Reemplaza la imagen del proceso (código y datos) por la del binario \texttt{path}. Si tiene éxito, no retorna. \\
\texttt{exit(exitstatus)} & Termina el proceso actual con el código \texttt{exitstatus}. \\
\texttt{wait(\&status)} & El padre se bloquea hasta que termine alguno de sus hijos. Retorna el estado de salida en \texttt{status}. \\
\texttt{getpid()} & Retorna el identificador del proceso. \\
\texttt{kill(pid)} & Envía una señal al proceso con identificador \texttt{pid}. \\
\texttt{sleep(n)} & Bloquea al proceso por \texttt{n} segundos. \\
\texttt{sbrk(n)} & Asigna \texttt{n} bytes adicionales al proceso. \\
\bottomrule
\end{tabular}
\end{center}

\begin{idea}
\texttt{fork()} y \texttt{exec()} suelen ir juntos: el shell hace
\texttt{fork()} para crear un hijo, y el hijo hace \texttt{exec()}
del programa que el usuario quiere ejecutar. Mientras tanto, el
padre se queda en \texttt{wait()} esperando a que el hijo termine.
Esta es la base del modelo de procesos de UNIX.
\end{idea}

\subsection{Estados de un proceso}

A lo largo de su vida, una tarea pasa por varios estados. Las Notas
identifican los siguientes:

\begin{itemize}[leftmargin=*,nosep]
    \item \textbf{NEW}: recién creada, todavía no ingresó al
        scheduler.
    \item \textbf{RUNNING}: actualmente ejecutando, tiene asignada
        una CPU.
    \item \textbf{RUNNABLE} (o \textbf{READY}): lista para ejecutar,
        esperando que el kernel le asigne una CPU.
    \item \textbf{SLEEPING} (o \textbf{WAITING}): bloqueada hasta la
        ocurrencia de un evento que la despertará (por ejemplo, la
        finalización de una operación de I/O).
    \item \textbf{ZOMBIE}: terminó (\texttt{exit()}) pero el padre
        todavía no ejecutó \texttt{wait()} para recoger su estado de
        salida.
    \item \textbf{TERMINATED}: una vez que el padre hace
        \texttt{wait()}, el descriptor se libera.
\end{itemize}

\begin{center}
\begin{tikzpicture}[
    node distance=2.0cm and 1.6cm,
    every node/.style={draw, font=\ttfamily\small, minimum width=1.8cm, minimum height=0.55cm},
    every edge/.style={draw,->,>=Stealth},
    edgelbl/.style={font=\scriptsize\sffamily, fill=white, inner sep=1pt}
]
    \node (new)   {NEW};
    \node (rdy)   [right=of new]   {RUNNABLE};
    \node (run)   [right=of rdy]   {RUNNING};
    \node (wait)  [below=of rdy]   {WAITING};
    \node (zomb)  [above=1.2cm of run]   {ZOMBIE};
    \node (term)  [above=of zomb]  {TERMINATED};

    \draw[->,>=Stealth] (new) -- node[edgelbl]{fork()} (rdy);
    \draw[->,>=Stealth] (rdy.north east) to[bend left=10] node[edgelbl]{scheduler()} (run.north west);
    \draw[->,>=Stealth] (run.south west) to[bend left=10] node[edgelbl]{yield()} (rdy.south east);
    \draw[->,>=Stealth] (run) -- node[edgelbl,sloped]{syscall()/suspend()} (wait);
    \draw[->,>=Stealth] (wait) -- node[edgelbl,sloped]{event/wakeup()} (rdy);
    \draw[->,>=Stealth] (run) -- node[edgelbl]{exit()} (zomb);
    \draw[->,>=Stealth] (zomb) -- node[edgelbl]{wait() del padre} (term);
\end{tikzpicture}
\\
\small\emph{Figura 3: Cambios de estado de un proceso.}
\end{center}

\paragraph{ZOMBIE y huérfanos.} Un proceso que terminó pero todavía
existe en la tabla del kernel se llama \emph{zombie}. Esto sucede
porque el SO debe preservar el código de salida hasta que el padre lo
recoja con \texttt{wait()}. Si el padre termina antes que el hijo, el
hijo queda \emph{huérfano} y es \emph{adoptado} por \texttt{init}
(PID~1), que se encarga de recogerlo cuando termine.

\subsection{Manejo de tareas en el kernel}

El kernel organiza las tareas en estructuras de datos. Las tareas en
estado \texttt{RUNNABLE} forman una cola (o varias colas, según el
algoritmo de scheduling). Las tareas en \texttt{SLEEPING} se encolan
en \emph{wait queues} asociadas al evento o recurso por el que
esperan: cuando llega el evento, el kernel desencola y despierta a la
tarea correcta.

\begin{idea}
Pensemos en una sala de espera del médico. La sala de listos
(\texttt{RUNNABLE}) tiene a todos los pacientes que pueden ser
atendidos. Las salas privadas (\texttt{wait queues}) tienen pacientes
esperando un evento puntual: un análisis, una llamada, un resultado.
Cuando el evento ocurre, el kernel los manda de vuelta a la sala de
listos.
\end{idea}

¿Qué pasa si el scheduler no encuentra ningún proceso \texttt{RUNNABLE}?
Hay dos alternativas:

\begin{enumerate}[leftmargin=*,nosep]
    \item Quedarse en un ciclo de busca activa.
    \item Tener un proceso \emph{idle} creado al inicio que ejecuta
        \texttt{while (true) yield();}. Al obtener la CPU la libera
        inmediatamente, asegurando que siempre hay un \texttt{RUNNABLE}
        que elegir.
\end{enumerate}

En cualquiera de los dos casos, este es el momento ideal para que el
SO contabilice la \emph{idleness} y dispare ahorro de energía o
trabajo diferido (procesar paquetes de red, por ejemplo) en
\emph{kernel threads}.

% =====================================================================
\section{Modos de ejecución, traps e interrupciones}
% =====================================================================

\subsection{Modo usuario y modo supervisor}

\begin{definicion}
Un proceso de usuario se ejecuta con la \textbf{CPU en modo usuario}.
En este modo la CPU no puede ejecutar ciertas instrucciones
privilegiadas (como deshabilitar interrupciones) y solo puede acceder
a su propio espacio de memoria. El kernel se ejecuta en \textbf{modo
kernel} (o \emph{supervisor}), un modo de mayor privilegio sin
restricciones.
\end{definicion}

Esta separación es la base de la protección: si un proceso pudiera
deshabilitar interrupciones o tocar la memoria del kernel, podría
hacer caer al sistema o monopolizar la CPU.

\subsection{Interrupciones, excepciones y traps}

Tres mecanismos llevan la CPU desde modo usuario a modo kernel:

\begin{itemize}[leftmargin=*,nosep]
    \item \textbf{Interrupciones}: las generan los dispositivos
        (teclado, discos, red) para notificar eventos
        (``terminé de leer'', ``llegó un paquete'').
    \item \textbf{Excepciones}: las genera la CPU al intentar ejecutar
        una instrucción inválida (división por cero, acceso a memoria
        ilegal).
    \item \textbf{Traps}: las dispara intencionalmente un programa de
        usuario para invocar al kernel; es la base de las
        \emph{llamadas al sistema}.
\end{itemize}

El SO instala, durante su inicialización, un \textbf{vector de
interrupciones}: una tabla en memoria del kernel donde cada entrada
apunta a una rutina (\emph{interrupt service routine} o ISR)
encargada de manejar el evento correspondiente.

\begin{center}
\begin{tikzpicture}[every node/.style={font=\ttfamily\small}]
    \node[draw, minimum width=2.4cm, minimum height=0.6cm,fill=yellow!10] (stvec) {stvec};
    \node[draw=none, above=0.1cm of stvec] {special cpu register};

    \matrix[draw, column sep=0pt, row sep=0pt, right=1.6cm of stvec, nodes={draw, minimum width=1.6cm, minimum height=0.6cm}] (vec) {
        \node (e0) {0}; \\
        \node (e1) {1}; \\
        \node (e2) {...}; \\
    };
    \node[draw=none, above=0.1cm of vec.north] {interrupts vector};

    \node[draw, right=2cm of e0,fill=blue!8!white] (isr0) {isr\_0: ... iret};
    \node[draw, right=2cm of e1,fill=blue!8!white] (isr1) {isr\_1: ... iret};

    \draw[->,>=Stealth] (stvec.east) -- (vec.west);
    \draw[->,>=Stealth] (e0.east) -- (isr0.west);
    \draw[->,>=Stealth] (e1.east) -- (isr1.west);
\end{tikzpicture}
\\
\small\emph{Figura 4: Vector de interrupciones.}
\end{center}

\subsection{Pasos del hardware ante una interrupción}

\begin{enumerate}[leftmargin=*,nosep]
    \item Salva el \texttt{pc} actual en algún registro o en memoria.
    \item Cambia el modo de ejecución a supervisor.
    \item Salta a la rutina correspondiente según el número de
        interrupción (índice del vector).
    \item Al retornar (instrucción \texttt{iret}), restaura el modo
        previo y el \texttt{pc} salvado.
\end{enumerate}

En arquitecturas como RISC-V, todas las interrupciones disparan una
única ISR cuya dirección está en \texttt{stvec}; el código del
manejador identifica el motivo leyendo \texttt{scause} y despacha al
caso correspondiente.

\begin{importante}
El kernel toma el control en cada interrupción, excepción o trap.
Esa es la única manera de que las funciones del kernel se invoquen
mientras un proceso de usuario está corriendo.
\end{importante}

\subsection{Implementación de syscalls}

Las syscalls se implementan en dos partes:

\begin{itemize}[leftmargin=*]
    \item \textbf{Lado usuario}: típicamente como funciones en la
        biblioteca estándar. Cada función pone el \emph{número de
        syscall} en un registro designado de la CPU (\texttt{a7} en
        RISC-V) y ejecuta una instrucción de trap (\texttt{int} en
        x86 o \texttt{ecall} en RISC-V).
    \item \textbf{Lado kernel}: el manejador de trap (ISR) reconoce
        que se trata de una syscall, lee el número y los argumentos
        del trapframe, y llama a la función correspondiente del
        kernel. El valor de retorno se escribe en el trapframe para
        que, al volver a usuario, esté en el registro adecuado.
\end{itemize}

Por convención POSIX, los valores negativos retornados por una
syscall (típicamente \texttt{-1}) son códigos de error.

% =====================================================================
\section{Cambios de contexto}
% =====================================================================

\subsection{El problema}

Si una sola CPU debe ejecutar muchas tareas, el SO necesita poder
\emph{intercambiar} cuál de ellas tiene la CPU en cada momento. Para
hacerlo correctamente, debe ser capaz de pausar a una tarea de modo
que, más adelante, pueda retomarla \emph{exactamente} donde la dejó.

\begin{definicion}
El \textbf{saved context} de una tarea es la estructura que guarda el
estado de la CPU desde que la tarea abandonó el procesador, para
poder continuar su ejecución más adelante. Típicamente contiene el
\texttt{pc}, el \texttt{sp} y posiblemente otros registros (los
\emph{callee-saved}).
\end{definicion}

El descriptor de cada tarea contiene su \texttt{tid}, su estado, sus
recursos y su \emph{saved context}.

\subsection{Quantum y preemption}

\begin{definicion}
El \textbf{quantum} (o \emph{time slice}) es el período de tiempo
máximo de uso continuo de la CPU asignado a una tarea. Cuando una
tarea agota su quantum, el kernel le \emph{quita} la CPU
(\emph{preemption}) y se la asigna a otra \texttt{RUNNABLE}.
\end{definicion}

\begin{idea}
El quantum hace que las tareas usen la CPU \emph{a ráfagas}, dando la
ilusión de que cada una tuviera su propia CPU. Si el quantum es muy
corto, hay demasiados context switches y el sistema gasta mucho
tiempo cambiando de tarea (overhead). Si es muy largo, las tareas
interactivas se sienten lentas. El equilibrio es un parámetro de
diseño del scheduler.
\end{idea}

\subsection{Operaciones \texttt{yield()} y \texttt{scheduler()}}

El kernel define dos funciones para implementar este mecanismo:

\begin{itemize}[leftmargin=*,nosep]
    \item \texttt{yield()}: el thread corriente \emph{abandona} la CPU.
        Marca su estado como \texttt{RUNNABLE} y dispara un context
        switch al \emph{scheduler}.
    \item \texttt{scheduler()}: selecciona el próximo thread a darle
        la CPU.
\end{itemize}

Esquema simplificado:

\begin{lstlisting}[style=cstyle]
struct task* current; // current task descriptor

void yield() {
    current->state = RUNNABLE;
    switch_context(&current->ctx, &scheduler_ctx);
}

void scheduler() {
    while (1) {
        current = NULL;
        struct task* task = select_runnable_task();  // scheduling algorithm
        if (task) {
            context_switch(&scheduler->ctx, &(task->ctx));
        }
    }
}
\end{lstlisting}

\subsection{El timer y el preemtive multitasking}

Un thread que no invoca \texttt{yield()} monopolizaría la CPU.
Para evitarlo, el kernel programa el \emph{timer} (reloj) para que
genere interrupciones periódicas. En cada tick, la ISR del timer
puede decidir invocar \texttt{yield()} si el quantum se agotó.

\begin{definicion}
\textbf{Preemptive multitasking} (o \emph{tiempo compartido}): forma
de planificación en la que el kernel puede recuperar la CPU de una
tarea sin la cooperación de ésta, gracias a las interrupciones del
timer.
\end{definicion}

\subsection{Pasos de un context switch}

Cuando el SO decide quitar la CPU a una tarea y asignársela a otra:

\begin{enumerate}[leftmargin=*,nosep]
    \item Salvar los registros de la CPU correspondientes en el
        contexto del descriptor de la tarea actual.
    \item Cambiar al stack de la nueva tarea.
    \item Restaurar los registros de la nueva tarea desde su contexto.
\end{enumerate}

Implementación típica de \texttt{context\_switch} (escrita en
assembly):

\begin{lstlisting}[style=cstyle]
context_switch(struct context *current, struct context *new) {
    save_ctx(current);   // copia pc y sp del HW al contexto current
    restore_ctx(new);    // copia pc y sp del contexto new al HW
}
\end{lstlisting}

\subsection{Esquema mínimo de un kernel}

Las Notas dan un esquema simplificado de cómo se ven las funciones de
manejo de tareas en el kernel. El \emph{trap handler} central despacha
según el motivo:

\begin{lstlisting}[style=cstyle]
trap(trap_number) {
    push_cpu_registers();
    switch (trap_number) {
        case TIMER:
            ticks++;
            set_next_timer_interrupt();
            if (current_task_used_quantum())
                yield();   // preempt CPU
            break;
        case SYSCALL:
            syscall();
            break;
        case DISK_IRQ:
            wakeup(&disk_wait_queue);
            break;
        // ...
    }
    restore_cpu_registers();
    return_from_trap();
}
\end{lstlisting}

Y las primitivas \texttt{suspend}/\texttt{wakeup} para bloqueo y
desbloqueo:

\begin{lstlisting}[style=cstyle]
suspend(struct wait_queue *wq) {
    current->state = WAITING;
    enqueue(wq, current);
    yield();
}

wakeup(struct wait_queue *wq) {
    struct task *task = dequeue(wq);
    if (task) task->state = RUNNABLE;
}
\end{lstlisting}

% =====================================================================
\section{Procesos zombie y el primer proceso del sistema}
% =====================================================================

\subsection{El par \texttt{exit}/\texttt{wait}}

\begin{itemize}[leftmargin=*,nosep]
    \item \texttt{exit(exitcode)}: finaliza el proceso corriente con
        código de salida \texttt{exitcode}.
    \item \texttt{wait(\&exitcode)}: el proceso corriente se bloquea
        hasta que finalice un hijo. El código de salida se copia en
        el parámetro \texttt{exitcode}.
\end{itemize}

Esto ofrece un mecanismo básico de sincronización y comunicación
entre procesos: el padre puede saber cómo terminó el hijo. Por
convención el código de salida $0$ significa terminación normal;
cualquier valor distinto, anormal.

\subsection{El primer proceso del sistema}

En sistemas tipo UNIX, el primer proceso lanzado es \texttt{init}
(PID~1). Construido manualmente por el kernel al final del boot, su
imagen inicial es \texttt{initcode}, que invoca \texttt{exec("init")}.
A partir de ahí, init dispara los procesos del sistema:

\begin{lstlisting}[style=textstyle]
initcode  pid=1   // hardcoded en el kernel
   | exec("init")
init      pid=1
   | fork(); child: exec("tty")
tty       pid=2
   | exec("login")
login     pid=3
   | exec("sh")
shell     pid=4
\end{lstlisting}

Esto hace que el sistema mantenga un \emph{árbol de procesos} con la
relación padre/hijo, donde init es la raíz. Cuando un padre muere
antes que sus hijos, los hijos son adoptados por init, que se
encargará del \texttt{wait()} correspondiente.

% =====================================================================
\section{Concurrencia}
% =====================================================================

\subsection{Definiciones y por qué importa}

\begin{definicion}
\textbf{Concurrencia}: mecanismos en un sistema que posibilitan la
ejecución de múltiples hilos de ejecución por medio de ejecución
simultánea (\emph{paralelismo}) o tiempo compartido (\emph{time
sharing}) con cambios de contexto. Estos mecanismos pueden generar
diferentes secuencias de operaciones de manera no determinística.
\end{definicion}

Aún con un único hilo de ejecución secuencial, el mecanismo de
interrupciones o señales genera diferentes secuencias de ejecución
(\emph{intercalados} o \emph{interleavings}). Por ejemplo:

\begin{lstlisting}[style=cstyle]
int v = 1;

// disparado por una senal
signal_handler() {
    v += 1;          // (1)
}

// thread principal
v *= 2;              // (2)
print(v);            // (3)
\end{lstlisting}

Hay tres interleavings posibles según cuándo llegue la señal:

\begin{enumerate}[leftmargin=*,nosep]
    \item Sin señal: $[(2),(3)]$ imprime \texttt{2}.
    \item Señal entre (2) y (3): $[(2),(1),(3)]$ imprime \texttt{3}.
    \item Señal antes de (2): $[(1),(2),(3)]$ imprime \texttt{4}.
\end{enumerate}

Tres salidas distintas para el mismo programa: ese es exactamente el
problema de la concurrencia.

\subsection{Paralelismo vs.\ tiempo compartido}

\begin{itemize}[leftmargin=*,nosep]
    \item \textbf{Paralelismo}: dos o más hilos se ejecutan
        \emph{realmente} a la vez en CPUs distintas. Es el caso de
        los sistemas multicore o multiprocesador.
    \item \textbf{Tiempo compartido}: una única CPU se va alternando
        entre varios hilos vía cambios de contexto. Solo uno corre
        en cada instante, pero al usuario le parece que todos
        progresan.
\end{itemize}

En ambos casos hay concurrencia: hilos concurrentes son aquellos cuya
relación de orden no está garantizada por el programa.

\subsection{Recursos compartidos, regiones críticas y race conditions}

Dos o más tareas \emph{cooperativas} usan algún recurso compartido:
una variable, un archivo, un dispositivo. En procesos con múltiples
threads, lo más común son las variables compartidas (\emph{shared
memory}). En procesos separados, hay que pasar mensajes o usar
shared memory explícita.

\begin{definicion}
La \textbf{región crítica} es la porción de código que accede a un
recurso compartido.
\end{definicion}

\begin{definicion}
Una \textbf{condición de carrera} (\emph{race condition}) surge cuando
el comportamiento del sistema depende de la secuencia (\emph{timing}
o \emph{interleaving}) de la ocurrencia de eventos externos que
pueden generar flujos de ejecución no deseados.
\end{definicion}

\begin{importante}
La concurrencia genera trazas de ejecución que no están bajo el
control del programador, sino del scheduler (\emph{no determinismo}).
En un sistema correcto, este no determinismo no debería alterar los
objetivos de cómputo del programa. La existencia de una condición de
carrera \emph{hace observable} ese no determinismo.
\end{importante}

Para impedir trazas no deseadas hace falta usar \textbf{mecanismos de
sincronización}.

% =====================================================================
\section{Locks y exclusión mutua}
% =====================================================================

\subsection{Definición}

\begin{definicion}
Un \textbf{lock} (o \emph{mutex}) se asocia a un recurso para
sincronizar su acceso de manera exclusiva. Provee \textbf{exclusión
mutua}: cuando un thread adquirió el lock, los demás que lo
requieran deberán esperar hasta que el primero lo libere.
\end{definicion}

Un lock típicamente expone dos operaciones: \texttt{acquire(lock)} y
\texttt{release(lock)}. El patrón estándar es:

\begin{lstlisting}[style=cstyle]
lock l = create_lock("recurso r");
// ...
acquire(&l);       // espera a que el recurso este libre
use_resource(r);   // region critica
release(&l);       // libera
\end{lstlisting}

\paragraph{Costos.} El uso de locks tiene consecuencias adversas:
las operaciones se \emph{secuencializan} y producen
\emph{contención} en los hilos que compiten por el recurso. Hay que
minimizar la duración de la región crítica para reducir la pérdida
de paralelismo.

\subsection{Spinlocks: el primer intento}

Un \emph{spinlock} hace que cada thread espere usando la CPU en un
ciclo (\emph{espera ocupada}). Una primera implementación:

\begin{lstlisting}[style=cstyle]
void acquire(int *lk) {
    while (*lk)
        ;            // (1) busy wait
    *lk = 1;        // (2)
}

void release(int *lk) {
    *lk = 0;
}
\end{lstlisting}

\paragraph{¿Por qué falla?} Supongamos que el thread~1 ejecuta
\texttt{acquire()} y está en la línea~(1) con \texttt{*lk == 0}.
Antes de la línea~(2) ocurre una interrupción de reloj y el scheduler
pasa al thread~2. El thread~2 también invoca \texttt{acquire()}, ve
\texttt{*lk == 0} y avanza a la línea~(2), entrando a la región
crítica. Otra interrupción y el control vuelve al thread~1, que
también ejecuta la línea~(2) y entra a la región crítica. \textbf{No
se pudo garantizar exclusión mutua}: el problema es que el lock en
sí mismo es un recurso compartido y las operaciones \texttt{acquire}
y \texttt{release} deberían ser \emph{atómicas} (sin interrupciones
en el medio).

\subsection{Solución 1: deshabilitar interrupciones}

En un monoprocesador podemos deshabilitar interrupciones en
\texttt{acquire()}:

\begin{lstlisting}[style=cstyle]
void acquire(int *lk) {
    cli();              // disable interrupts
    while (*lk) ;
    *lk = 1;
}

void release(int *lk) {
    *lk = 0;
    sti();              // enable interrupts
}
\end{lstlisting}

Funciona en monoprocesadores. \emph{No funciona en multiprocesadores}:
\texttt{cli()} solo deshabilita interrupciones en la CPU que la
ejecuta, no en las otras.

\subsection{Solución 2: instrucciones atómicas}

Las CPUs modernas multicore con SMP proveen instrucciones atómicas
indivisibles, implementadas mediante bloqueo del bus de memoria,
bloqueo de la línea de caché o, en monoprocesadores, deshabilitando
interrupciones. Las más comunes son:

\begin{itemize}[leftmargin=*]
    \item \textbf{Compare and swap (CAS)}:
        \texttt{cas(addr, old, new)}: si la variable en \texttt{addr}
        contiene \texttt{old}, la modifica con \texttt{new}.
    \item \textbf{Test and set (TAS)}:
        \texttt{old = tas(addr, new)}: asigna \texttt{new} y retorna
        el valor anterior.
    \item \textbf{Fetch and add (FAA)}:
        \texttt{faa(addr, value)}: suma \texttt{value} atómicamente.
\end{itemize}

GCC provee \texttt{\_\_sync\_lock\_test\_and\_set(\&lock, value)}
con implementación dependiente de la arquitectura. En RISC-V emite:

\begin{lstlisting}[style=textstyle]
li     a5, 1
amoswap.w a5, a5, (lock_address)
\end{lstlisting}

\subsection{Spinlocks en multiprocesadores}

Una implementación correcta de spinlock para SMP combina
deshabilitar interrupciones, instrucción atómica y \emph{barreras
de memoria}:

\begin{lstlisting}[style=cstyle]
void acquire(int *lk) {
    cli();
    while (__sync_lock_test_and_set(lk, 1) != 0)
        ;
    __sync_synchronize();   // memory barrier (fence)
}

void release(int *lk) {
    __sync_lock_release(lk); // *lk = 0;
    __sync_synchronize();
    sti();
}
\end{lstlisting}

\paragraph{Por qué la barrera.} En multiprocesadores con \emph{modelo
de memoria relajado} (lo habitual en x86, ARM, RISC-V), una escritura
hecha por una CPU puede no ser visible inmediatamente en otra (queda
en su caché L1). La barrera fuerza a la CPU a propagar la escritura
a memoria principal e invalidar las cachés ajenas, asegurando que el
cambio del lock sea visible en todos los cores.

\subsection{Sleep locks}

Un spinlock solo es útil para esperas \emph{cortas}: si la región
crítica tarda mucho, gastar CPU haciendo busy-wait es un desperdicio.
Cuando la espera puede ser larga, conviene \emph{suspender} al
thread y reactivarlo cuando la condición cambie. Esto es lo que hace
un \textbf{sleep lock}, similar a una variable de condición:

\begin{itemize}[leftmargin=*,nosep]
    \item \texttt{sleep(task* t, wait\_queue* q)}: pasa $t$ a estado
        \texttt{SLEEPING} y lo encola en $q$.
    \item \texttt{wakeup(wait\_queue* q)}: desencola una tarea y la
        pasa a \texttt{RUNNABLE}.
\end{itemize}

La implementación interna requiere de spinlocks porque la cola es a
su vez un recurso compartido y \texttt{sleep}/\texttt{wakeup} pueden
ejecutarse concurrentemente.

\subsection{Modularidad y problemas del uso de locks}

\begin{importante}
Los locks \textbf{no son modulares}. Un thread puede adquirir un
lock en una función y liberarlo en otra; los efectos colaterales del
uso de locks deben ser documentados explícitamente.
\end{importante}

El uso de locks es propenso a tres problemas clásicos:

\paragraph{\emph{Self-lock} o \emph{double lock}.} Si un thread
intenta adquirir un recurso que él mismo ya posee, queda bloqueado
para siempre. Solución: \emph{locks recursivos}.

\paragraph{\emph{Deadlock} (interbloqueo).} Dos hilos esperan
mutuamente un recurso del otro:

\begin{lstlisting}[style=cstyle]
f() {                      g() {
    acquire(&l1);              acquire(&l2);
    // ...                     // ...
    acquire(&l2);              acquire(&l1);   // (potencial deadlock)
    release(&l2);              release(&l1);
    release(&l1);              release(&l2);
}                          }
\end{lstlisting}

Si el thread~A ejecutando \texttt{f()} adquiere \texttt{l1} y el
thread~B ejecutando \texttt{g()} adquiere \texttt{l2}, ambos quedan
esperando indefinidamente. Las cuatro condiciones para que exista
posibilidad de deadlock son (\emph{condiciones de Coffman}):

\begin{enumerate}[leftmargin=*,nosep]
    \item \textbf{Exclusión mutua}: los recursos se adquieren de
        forma exclusiva.
    \item \textbf{Espera y retención}: un proceso retiene un recurso
        mientras espera por otro.
    \item \textbf{No quita} (\emph{preemption}): los recursos se
        liberan solo voluntariamente.
    \item \textbf{Espera circular}: existe una cadena
        $P_0, P_1, \dots, P_{n-1}, P_i$ tal que cada uno espera por un
        recurso adquirido por $P_{(i+1) \bmod n}$.
\end{enumerate}

La estrategia más habitual para prevenir deadlock es \emph{romper la
condición 4} estableciendo un orden total para la adquisición de
locks: si todos los módulos del sistema toman \texttt{l1} antes que
\texttt{l2}, jamás se forma una cadena circular.

\paragraph{\emph{Livelock}.} Dos o más hilos quedan en ciclos sin
progresar en alguna computación útil. No están bloqueados (siguen
ejecutando), pero ninguno avanza. Es típico de protocolos mal
diseñados de \emph{retry}.

% =====================================================================
\section{Solución de Peterson}
% =====================================================================

Edsger Dijkstra y Gary Peterson estudiaron en los años 70 cómo
resolver el problema de la región crítica usando \emph{solo} lectura
y escritura sobre memoria compartida (sin instrucciones atómicas del
hardware). El algoritmo de Peterson (1981) lo logra para dos threads:

\begin{lstlisting}[style=cstyle]
bool flag[2] = {false, false};
int  turn;

lock(i) {                      unlock(i) {
    flag[i] = true;                flag[i] = 0;
    int other = (i==0)? 1: 0;  }
    turn = other;
    while (flag[other] && turn == other)
        ;
}
\end{lstlisting}

La idea: \texttt{flag[i]} expresa la \emph{intención} del thread $i$
de ingresar a la sección crítica; \texttt{turn} indica qué thread
tendría la prioridad para ingresar en caso de empate. Anunciando el
interés \emph{antes} de ceder el turno, el algoritmo garantiza
exclusión mutua y progreso (la demostración semi-formal está en el
Ejercicio~10 del Práctico~3).

El algoritmo se generaliza a $N$ threads (\emph{algoritmo filter}),
pero su valor en la práctica moderna es limitado:

\begin{importante}
La solución de Peterson \textbf{no funciona en sistemas
multiprocesadores} sin barreras de memoria explícitas, debido a las
cachés y los modelos de memoria relajados. Por eso los sistemas
reales usan instrucciones atómicas del hardware (TAS, CAS, FAA) y
construyen mutexes/spinlocks sobre ellas.
\end{importante}

% =====================================================================
\section{Semáforos}
% =====================================================================

\subsection{Definición}

\begin{definicion}
Un \textbf{semáforo} (Dijkstra, 1963) representa una cantidad de
recursos disponibles. Es una primitiva de sincronización más general
que los locks. Su API tiene tres operaciones:
\begin{itemize}[leftmargin=*,nosep]
    \item \texttt{sem\_init(s, value)}: inicializa el valor del
        semáforo $s$.
    \item \texttt{sem\_wait(s)} (\emph{P}): si el valor de $s$ es
        cero, suspende al invocante; si no, decrementa.
    \item \texttt{sem\_signal(s)} (\emph{V}, también
        \texttt{sem\_post}): incrementa el valor de $s$ y despierta a
        algún thread bloqueado.
\end{itemize}
\end{definicion}

Un semáforo $n$-ario (o \emph{counting semaphore}) permite que hasta
$n$ threads accedan al recurso. Un semáforo binario (con $n = 1$) es
equivalente a un \textbf{mutex}: garantiza exclusión mutua.

Internamente un semáforo se implementa con un contador y un sleep
lock: si el contador llega a cero, los \texttt{wait()}s subsiguientes
ponen al thread en una wait queue; un \texttt{signal()} desencola a
uno y lo pasa a \texttt{RUNNABLE}.

\subsection{El productor-consumidor con semáforos}

El problema clásico: uno o más productores generan valores y los
ponen en un buffer acotado de capacidad $N$; uno o más consumidores
extraen valores del buffer y los procesan. Productores y consumidores
ejecutan concurrentemente.

\begin{lstlisting}[style=cstyle]
buffer buf[N];
semaphore mutex = 1, full = N, empty = 0;

void producer() {
    while (true) {
        value = gen_value();
        sem_wait(&full);
        add_to_buffer(value);
        sem_signal(&empty);
    }
}

void consumer() {
    while (true) {
        sem_wait(&empty);
        value = get_from_buffer();
        sem_signal(&full);
        process(value);
    }
}

void add_to_buffer(value) {
    sem_wait(&mutex);
    // agregar value al buffer
    sem_signal(&mutex);
}

item_type get_from_buffer() {
    sem_wait(&mutex);
    // extraer item
    sem_signal(&mutex);
}
\end{lstlisting}

\begin{idea}
Tres semáforos cumplen tres roles distintos:
\begin{itemize}[leftmargin=*,nosep]
    \item \texttt{full}/\texttt{empty}: \emph{condición} (cuántos
        slots libres u ocupados hay).
    \item \texttt{mutex}: \emph{exclusión mutua} sobre la
        modificación del buffer y los índices.
\end{itemize}
Las regiones críticas son lo más cortas posible para no serializar
todo el pipeline.
\end{idea}

% =====================================================================
\section{Variables de condición y monitores}
% =====================================================================

\subsection{Variables de condición}

\begin{definicion}
Una \textbf{variable de condición} es una primitiva de sincronización
que maneja una cola de tareas esperando una condición (estado
lógico). Comúnmente se usa con un lock asociado al recurso. Su
interfaz es:
\begin{itemize}[leftmargin=*,nosep]
    \item \texttt{wait(\&cond, \&lock)}: requiere que el invocante haya
        adquirido el lock. \emph{Atómicamente} libera el lock, agrega
        el thread a la cola y lo suspende.
    \item \texttt{signal(\&cond)}: invocada por un thread que cambió
        la condición. Despierta uno o todos los threads esperando.
\end{itemize}
\end{definicion}

Productor-consumidor usando variables de condición:

\begin{lstlisting}[style=cstyle]
condition cond; mutex lock; buffer buf;

producer() {                         consumer() {
    acquire(&lock);                      acquire(&lock);
    while (buffer_is_full())             while (buffer_is_empty())
        wait(&cond, &lock);                  wait(&cond, &lock);
    produce(&buf);                       consume(&buf);
    signal(&cond);                       signal(&cond);
    release(&lock);                      release(&lock);
}                                    }
\end{lstlisting}

Las operaciones \texttt{wait()}/\texttt{notify()} de Java
corresponden exactamente a este mecanismo.

\subsection{Monitores}

\begin{definicion}
Un \textbf{monitor} (Hansen y Hoare, 1974) es una construcción
sintáctica de un lenguaje de programación que encapsula datos y las
operaciones sobre ellos al estilo orientado a objetos. Los hilos
solo pueden manipular los datos por medio de las operaciones del
monitor. Es \emph{thread safe}: los invariantes están protegidos
ante concurrencia.
\end{definicion}

Implementación común: variables de condición + al menos un mutex que
asegura la exclusión mutua entre operaciones concurrentes.

\begin{lstlisting}[style=cstyle]
monitor PC {
    condition full, empty;
    data buffer[N];
    int count = 0;

    void put(data item) {
        if (count == N) wait(full);
        insert(buffer, item);
        if (++count == 1) signal(empty);
    }

    data get() {
        if (count == 0) wait(empty);
        data item = remove(buffer);
        if (--count == N-1) signal(full);
        return item;
    }
};

void producer() {
    while (true) { data item = gen_item(); PC.put(item); }
}
void consumer() {
    while (true) { data item = PC.get(); process(item); }
}
\end{lstlisting}

\begin{idea}
Los monitores son la abstracción \emph{modular} de la sincronización:
desde afuera, el productor y el consumidor solo ven \texttt{put()} y
\texttt{get()}. Toda la complejidad (mutex, variables de condición,
cuándo \texttt{wait}/\texttt{signal}) queda encapsulada y el código
cliente no se puede equivocar olvidándose un \texttt{release}.
\end{idea}

Java soporta monitores con la palabra clave \texttt{synchronized}:
una clase con métodos sincronizados es un monitor implícito. Cada
objeto Java tiene asociado un mutex y una variable de condición
disponibles vía \texttt{wait()}/\texttt{notify()}/\texttt{notifyAll()}.

% =====================================================================
\section{Mensajes y comunicación}
% =====================================================================

\subsection{Mensajes en lugar de memoria compartida}

Los mecanismos vistos hasta acá se basan en \emph{shared memory}.
Cuando los procesos no comparten memoria (corren en computadoras
distintas o el SO no les permite memoria compartida), se necesita
\emph{pasaje de mensajes}. El modelo se basa en dos primitivas:

\begin{itemize}[leftmargin=*,nosep]
    \item \texttt{send(destination, message)}: envía un mensaje a un
        proceso identificado.
    \item \texttt{receive(source, \&message)}: recibe un mensaje del
        emisor (puede ser \texttt{any}).
\end{itemize}

La comunicación puede ser:

\begin{itemize}[leftmargin=*,nosep]
    \item \textbf{Asincrónica}: hay un buffer de mensajes
        (\emph{mailbox}). El emisor no se bloquea; el receptor solo
        se bloquea si el mailbox está vacío.
    \item \textbf{Sincrónica} (\emph{rendezvous}): no hay buffer.
        Emisor y receptor deben coincidir en \texttt{send}/\texttt{receive}.
\end{itemize}

\subsection{Micro-kernels}

En SOs con diseño \emph{micro-kernel} como MINIX, los subsistemas
(gestor de procesos, memoria, sistemas de archivos) se implementan
como \emph{servers} que reciben requerimientos por mensajes:

\begin{lstlisting}[style=cstyle]
int main(void) {
    message m;
    init_subsystem();
    while (TRUE) {
        receive(ANY, &m);
        switch (m.type) {
            case REQUEST_1:
                // ...
                send(other_subsystem, m2);
                break;
            // ...
        }
    }
}
\end{lstlisting}

% =====================================================================
\section{Futuros, promesas y concurrencia sin locks}
% =====================================================================

\subsection{Futuros y promesas}

\begin{definicion}
Un \textbf{futuro} (o \emph{promesa}) es una abstracción que
representa un valor que será computado por algún otro componente de
ejecución concurrente. Se usan para implementar \emph{data-flow
variables}: variables que ante una lectura esperan a que se asigne su
valor. La operación de lectura actúa como mecanismo de sincronización.
\end{definicion}

\begin{itemize}[leftmargin=*,nosep]
    \item \texttt{future f = async g(args)}: invocación asincrónica
        de \texttt{g()} en un nuevo thread. El invocante continúa
        ejecutando.
    \item \texttt{value = f.get()}: obtiene el valor encapsulado en
        $f$, posiblemente bloqueándose hasta que el thread que evalúa
        \texttt{g()} finalice.
\end{itemize}

\subsection{Estructuras lock-free y wait-free}

\begin{definicion}
Una estructura de datos es \textbf{lock-free} si puede usarse
concurrentemente sin operaciones de bloqueo. Es \textbf{wait-free} si
cada thread que la usa puede completar sus operaciones en un número
acotado de pasos.
\end{definicion}

Un spinlock califica como lock-free pero no como wait-free. Las
estructuras lock-free generalmente requieren instrucciones atómicas
y tienen ventajas de rendimiento (no hay contención por locks) y
seguridad (si un thread finaliza con un lock adquirido, los demás no
quedan bloqueados). Ejemplo de \texttt{push} lock-free a una lista
enlazada usando CAS:

\begin{lstlisting}[style=cstyle]
push(list* head, T value) {
    list_node *n = create_node(value);
    n->next = head->first;
    while (!cas(&(head->first), n->next, n))
        ;
}
\end{lstlisting}

% =====================================================================
\section{IPC en UNIX}
% =====================================================================

Los sistemas tipo UNIX ofrecen una variedad de mecanismos para la
comunicación y sincronización entre procesos:

\begin{center}
\small
\begin{tabular}{@{}lp{8.5cm}@{}}
\toprule
Mecanismo & Llamadas al sistema \\
\midrule
File                & \texttt{open, read, write, close, ...} \\
Pipes               & \texttt{pipe, read, write, close, ...} \\
Named pipes         & \texttt{mkfifo, read, write, close, ...} \\
Signals             & \texttt{kill, signal, alarm, ...} \\
Message queues      & \texttt{msgget, msgrcv, msgctl, ...} \\
Semaphores          & \texttt{semget, semctl, semop, sem\_open, sem\_wait, sem\_post, ...} \\
Shared memory       & \texttt{shmget, shmat, shmctl, shm\_open, mmap, ...} \\
Memory mapped files & \texttt{mmap, munmap, ...} \\
Sockets             & \texttt{listen, connect, accept, send, recv, ...} \\
\bottomrule
\end{tabular}
\end{center}

Todos estos mecanismos permiten la comunicación entre procesos
locales, salvo los \emph{sockets}, que también permiten la
comunicación entre procesos en computadoras distintas (red).

% =====================================================================
\section{Resumen final}
% =====================================================================

\begin{enumerate}[leftmargin=*]
    \item Un \textbf{proceso} es una instancia de un programa en
        ejecución, con espacio de memoria propio. Un \textbf{thread}
        es la mínima unidad planificable: comparte memoria con otros
        threads del mismo proceso, pero tiene pila propia.
    \item El SO mantiene por cada tarea un \textbf{descriptor} (PCB)
        con su PID, estado, mapa de memoria, recursos y \emph{saved
        context}. Las tareas en \texttt{RUNNABLE} están en cola de
        listos; las \texttt{SLEEPING}, en wait queues asociadas al
        evento por el que esperan.
    \item Las transiciones de estado se disparan por syscalls,
        interrupciones (timer, dispositivos) y excepciones. El timer
        es el mecanismo central del \textbf{preemptive multitasking}:
        permite al kernel recuperar la CPU sin la cooperación del
        proceso.
    \item Un \textbf{cambio de contexto} guarda los registros del
        thread saliente y restaura los del entrante. La operación
        clave es \texttt{context\_switch(current, new)} y suele
        escribirse en assembly.
    \item La \textbf{concurrencia} es no determinismo controlado por
        el scheduler. Cuando varias tareas comparten un recurso, los
        \emph{interleavings} pueden producir
        \textbf{condiciones de carrera}, que se evitan protegiendo
        las regiones críticas con primitivas de sincronización.
    \item Hay una jerarquía de primitivas:
        \begin{itemize}[leftmargin=*,nosep]
            \item \emph{Spinlocks}: para regiones críticas cortas.
                Usan instrucciones atómicas y barreras de memoria en
                multiprocesadores.
            \item \emph{Sleep locks}: para regiones largas. Suspenden
                al thread y lo despiertan con \texttt{wakeup}.
            \item \emph{Semáforos}: generalización del lock; el
                contador modela cantidad de recursos.
            \item \emph{Variables de condición}: hilos esperan que
                otra tarea cambie un estado lógico.
            \item \emph{Monitores}: encapsulación modular de los
                anteriores en una clase \emph{thread safe}.
        \end{itemize}
    \item El uso de locks introduce problemas: \emph{contención},
        \emph{deadlock} (cuatro condiciones de Coffman),
        \emph{livelock}, \emph{self-lock}. La estrategia más habitual
        para prevenir deadlock es definir un orden total de
        adquisición.
    \item El algoritmo de \textbf{Peterson} demuestra que la
        exclusión mutua entre dos threads puede resolverse solo con
        lecturas y escrituras a memoria compartida, pero requiere
        atomicidad y consistencia secuencial; en la práctica los
        sistemas usan instrucciones atómicas del hardware.
    \item El \textbf{productor-consumidor} con tres semáforos
        (\texttt{full}, \texttt{empty}, \texttt{mutex}) es el
        ejemplo prototípico de combinar sincronización por condición
        y exclusión mutua.
    \item Cuando los procesos no comparten memoria, se comunican por
        \textbf{mensajes} (sincrónicos o asincrónicos). UNIX ofrece
        un catálogo amplio de mecanismos IPC: pipes, FIFOs, signals,
        message queues, semáforos, shared memory, sockets.
\end{enumerate}

\begin{importante}
\textbf{La idea unificadora:} la concurrencia es no determinismo
controlado, y el rol del SO es ofrecer abstracciones suficientes
(estados, scheduler, locks, semáforos, monitores, mensajes) para que
ese no determinismo no altere los objetivos de cómputo de los
programas.
\end{importante}

\end{document}
